\chapter{Abstract}

Due to recent technological advances in both Hardware and Software-Integration, like the release of the Apple Vision Pro\footnote{https://www.apple.com/apple-vision-pro/},  the Meta Ray-Ban Display\footnote{https://www.meta.com/en-gb/blog/meta-ray-ban-display-ai-glasses-connect-2025/} and the gradual release of further developer SDK-features on the Meta Horizon-Platform, like the new Passthrough API\footnote{https://developers.meta.com/horizon/documentation/unity/unity-passthrough/}, Augmented Reality Applications are gradually becoming more popular.
\\
In Augmented Reality the real environment is shown to the user, enhanced through additional virtual objects.

For Use-Cases like immersive analytics, in which the users focus should be on the visual data analysis, this visual clutter may prove disadvantageous and reduce the benefits provided by the interactive and flexible approach of analyzing data in an immersive three-Dimensional environment.
\\\\
This bachelors thesis focuses on the analysis of the different methods and algorithms for the concealment of the clutter to provide a clean and non-distracting environment for the digital object to be displayed on. 
\\
To accomplish this, a solution to select planes on which objects are detected to subsequently be hidden through inpainting will be implemented for the Unity3D Game-Engine. Through the usage of this solution, several methods and algorithms will be evaluated in terms of both quality and performance to provide believable and uncluttered environments, whilst at the same time minimizing the use of computational resources and time.
\\ 
These findings are compared to each other to find both the advantages and disadvantages each of the analyzed approaches delivers.